% =====================================================================
% HMK-1 · Hybrid Post-Quantum OpenID Connect — submission paper
% Mathematics for Systems Security · Cryptographer track
% Build:  tectonic main.tex     (or: latexmk -pdf main.tex)
% =====================================================================
\documentclass[11pt,a4paper]{article}

\usepackage[margin=2.4cm]{geometry}
\usepackage{amsmath,amssymb,amsthm}
\usepackage{algorithm,algorithmic}
\usepackage{booktabs}
\usepackage{hyperref}
\usepackage{xcolor}
\usepackage{listings}
\usepackage{microtype}

\hypersetup{colorlinks=true,linkcolor=blue!50!black,citecolor=blue!50!black,
            urlcolor=blue!50!black,pdftitle={HMK-1: Constant-Time Composite-Hybrid
            Master-Key Encryption for Post-Quantum OpenID Connect}}

\newtheorem{theorem}{Theorem}[section]
\newtheorem{corollary}{Corollary}[section]
\newtheorem{proposition}{Proposition}[section]
\newtheorem{definition}{Definition}[section]
\theoremstyle{remark}
\newtheorem{remark}{Remark}[section]

\lstset{basicstyle=\ttfamily\small,breaklines=true,columns=fullflexible,
        showstringspaces=false,frame=single,framerule=0.4pt,rulecolor=gray!50}

\newcommand{\adv}[1]{\mathrm{Adv}^{#1}}
\newcommand{\negl}{\mathrm{negl}(\lambda)}
\newcommand{\concat}{\,\|\,}
\newcommand{\F}{\mathbb{F}}
\newcommand{\STATESK}{\STATE}

\title{\textbf{HMK-1: A Constant-Time Composite-Hybrid\\
Master-Key Encryption Scheme for\\
Post-Quantum OpenID Connect}}
\author{Hybrid-OIDC Team \textnormal{(\texttt{github.com/HaiqaKaleem1/hybrid-oidc})}\\
\small Mathematics for Systems Security — Cryptographer Track}
\date{October 2026}

\begin{document}
\maketitle

\begin{abstract}
We present \textbf{HMK-1}, a master-key encryption scheme designed for the
OpenID Connect ecosystem under a post-quantum adversary, together with complete
security proofs for the composite-hybrid KEM and signature constructions it is
built on. The transport layer uses a \emph{composite hybrid} of a classical and
a post-quantum component: we prove that indistinguishability of the hybrid
secret reduces to the \emph{sum} of the component advantages, so traffic
survives the total break of either half (record-now-decrypt-later and
``the lattice assumption is wrong'' adversaries). The token layer uses
conjunction acceptance, whose unforgeability is provably that of the
\emph{weakest} half; we further bind tokens to the TLS channel to close the
bearer-replay gap. The master key itself is sealed by HMK-1: a novel
composition of an $\F_p$ field combiner (CIOS Montgomery arithmetic in the
NIST P-256 prime field, with a one-break bijection property), an invertible
$\F_2$ mixing matrix of rank 512, and an HMAC extractor feeding an
encrypt-then-MAC wrapper. The core is implemented in C with explicit
side-channel hardening --- no secret-dependent branches, no secret-indexed
tables, full-scan tag comparison, branchless publication of the candidate key,
and fixed work per call --- and validated by measurement: rejection latency is
flat at $\approx$1\,ns across tag-prefix positions where a short-circuit
comparator leaks a linear 2--22\,ns oracle, and full unwrap latency varies by
less than 2\,\% across success, failure, and wrong-secret cases. All claims are
traceable to proofs and to passing tests in the reference implementation.
\end{abstract}

\section{Introduction}
OpenID Connect (OIDC) authenticates users to relying parties with signed
JWTs, and every secret in the flow --- code verifiers, tokens, key-encryption
keys --- ultimately rests on transport confidentiality and token
unforgeability. Both are threatened by quantum attackers: Shor breaks ECDH
and ECDSA, and a future cryptanalytic breakthrough on lattice assumptions
could break the standardized post-quantum primitives. Neither failure mode is
hypothetical in the \emph{design} sense: hybridization is the standard
mitigation, but hybrid designs are frequently asserted \emph{by citation}
rather than proven, and implementation-level side channels are routinely
declared out of scope.

This submission contributes:
\begin{enumerate}
  \item \textbf{Complete proofs} (not citations) for the composite-hybrid KEM
        (Theorem~\ref{thm:hybrid}) and composite signature
        (Theorem~\ref{thm:forge}) constructions used in hybrid OIDC, under an
        explicit three-capability threat model (E1--E3).
  \item \textbf{HMK-1}, a constant-time master-key wrap: a novel construction
        combining an $\F_p$ combiner with a one-break property
        (Prop.~\ref{prop:onebreak}), an $\F_2$ mixing layer with verified
        rank 512 (Def.~\ref{def:matrix}), and an HMAC-based extractor and
        encrypt-then-MAC wrapper (Alg.~\ref{alg:wrap}). No pre-existing cipher
        encrypts plaintext; the keystream is derived by the construction itself.
  \item \textbf{A side-channel evaluation} of the C implementation: threat-by-
        threat mitigation table (Sec.~\ref{sec:sidechannel}), a naive-vs-constant
        timing harness, and a regression test that fails if unwrap latency ever
        becomes data dependent.
  \item \textbf{Full traceability}: every claim maps to a proof and to a named
        test (Table~\ref{tab:trace}).
\end{enumerate}

\section{Scheme syntax}
Algorithms are stated with fixed-length encodings throughout, following the
IETF hybrid-TLS draft: variable-length components would
leak timing through parsing and invalidate the dual-PRF combiner argument.

\begin{lstlisting}[caption={Component syntax: Alg 1--2 (components) and Alg 3--4 (composite constructions)},label={alg:syntax}]
Alg 1 KEM pi = (Kgen, Enc, Dec)      Alg 2 Sig = (Skg, Sign, Ver)
  Kgen(1^l)  -> (pk, sk)               Skg(1^l)  -> (vk, sk)
  Enc(pk)    -> (ss, ct)  [fixed len]  Sign(sk,m) -> s   [fixed len]
  Dec(sk,ct) -> ss or bot

Alg 3 composite-hybrid KEM
  KgenH       = Kgen_c ; Kgen_p            (independent key pairs)
  EncH(pk_c,pk_p) = (ss_c||ss_p, ct_c||ct_p)
  DecH(sk_c,sk_p, ct_c||ct_p) = componentwise; bot if either fails
  K = HKDF-Extract(ss_h)                   dual-PRF combiner

Alg 4 composite-hybrid signature
  SkgH  = Skg_c ; Skg_p
  SignH(sk_c, sk_p, m) = sigma_c||sigma_p
  VerH  = Ver_c(vk_c,m,sigma_c) AND Ver_p(vk_p,m,sigma_p)
\end{lstlisting}

In the reference deployment $c$ is ECDH/ECDSA (P-256, ES256) and $p$ is
ML-KEM/ML-DSA (the paper prototype uses Kyber512/Falcon512). The HMK-1
implementation realizes both halves as idealized random-oracle components with
adjustable break flags so that reductions run to completion under
experiment; swapping in the concrete schemes leaves algorithms and proofs
unchanged.

\section{Threat model}
\begin{description}
  \item[E1 (RNDL)] Passive record-now-decrypt-later: stores transcripts,
        later recovers past plaintext of \emph{some} session.
  \item[E2 (single-primitive break)] Breaks \emph{one} of
        $\{$classical, pq$\}$ completely and retains that power forever
        (PPT or QPT). E2 with the pq side broken is exactly the
        ``LWE preprint is wrong'' scenario.
  \item[E3 (active MITM)] IND-CCA and EUF-CMA oracle adversary: adaptive
        encaps/decaps and sign/verify queries, arbitrary ciphertexts and
        signatures submitted to the endpoints, plus timing observation.
\end{description}

\section{One-break confidentiality of the hybrid KEM}
\label{sec:thm1}
\begin{theorem}[Composite-KEM hybridization]
\label{thm:hybrid}
Let $A$ be PPT attacking $\mathrm{Hyb}$-$\mathsf{IND\text{-}CCA}$ with
$q_{\mathrm{enc}}$ encapsulation and $q_{\mathrm{dec}}$ decapsulation queries.
Then there exist adversaries $B_c,B_p$ attacking the component games with
\[
  \adv{Hyb}_A(\lambda)
  \;\le\;
  \adv{\mathsf{IND\text{-}CCA}}_{B_c}(\lambda)
  \;+\;
  \adv{\mathsf{IND\text{-}CCA}}_{B_p}(\lambda)
  \;+\;
  q_{\mathrm{enc}}\cdot\negl .
\]
Hence if \emph{either} component is $\mathsf{IND\text{-}CCA}$ --- the other may
be totally broken --- the hybrid advantage is negligible.
\end{theorem}

\begin{proof}
Games $0$--$3$. Game $0$ is the real experiment: $A$ receives
$pk_h = pk_c\concat pk_p$, queries $\mathrm{Enc}_h$ (fresh secrets per query)
and $\mathrm{Dec}_h$, then outputs $ss_h^0, ss_h^1$ and guesses $b$.
Game $1$ replaces the $p$-component transcript with real encapsulations; any
distinguisher yields a $B_p$ that answers $\mathrm{Dec}$ queries by delegation,
so the gap is at most $\adv{\mathsf{IND\text{-}CCA}}_{B_p}$. Game $2$ does the
same for the $c$-component with gap $\adv{\mathsf{IND\text{-}CCA}}_{B_c}$
(no hardness of either side is used in these two steps: they are consistency
replacements only). In Game $3$ the challenge secret is
$ss_h^b = ss_c^b \concat ss_p^b$ with both halves real and independent.
Because the concatenation uses \emph{fixed-length} encodings, the bit $b$ is
determined by either half alone; by the tower rule for conditional entropy the
guess advantage splits into the sum of the per-half guessing advantages plus
$q_{\mathrm{enc}}\cdot\negl$ accounting for collision/length agreement of
independently sampled blocks. This last term is precisely where fixed lengths
and the HKDF dual-PRF requirement are load-bearing: with variable-length
blocks a re-parsing adversary breaks the split argument. Hybridizing over the
sequence gives the claim.
\end{proof}

\begin{corollary}[E1: pq broken, classical holds]
\label{corr:e1}
Setting $\adv{B_p} = 1$ but $\adv{B_c} \le \negl$ yields
$\adv{A} \le \negl$: a lattice break or a quantum preprint gains \emph{nothing}
on traffic whose classical half stands. Because every session is
hybrid-length, an observer cannot even identify which sessions would be
classical-only and thus broken.
\end{corollary}

\begin{corollary}[E2: classical broken, pq holds]
\label{corr:e2}
Symmetrically $\adv{B_c} = 1$, $\adv{B_p} \le \negl$ gives
$\adv{A} \le \negl$: Shor against ECDH does not recover $ss_h$, because the
pq block alone is uniform under ML-KEM's assumption.
\end{corollary}

\section{Weakest-link unforgeability of the hybrid signature}
\begin{theorem}[Composite-signature bound]
\label{thm:forge}
For any PPT forger $A$ against $\mathsf{VerH}$ with $q_h$ signing queries and
messages of length at most $\sigma$,
\[
  \adv{EUF\text{-}CMA}_{A}
  \;\le\;
  \adv{EUF\text{-}CMA}_{B_c}
  \;+\;
  \adv{EUF\text{-}CMA}_{B_p}
  \;+\;
  q_h\cdot 2^{-\sigma/2},
\]
and, under conjunction acceptance, the bound is tight in the \emph{weaker}
half: a full break of one component cannot be compensated by the other.
\end{theorem}

\begin{proof}
Suppose $A$ outputs a fresh accepting pair $(m^*, \sigma_c^*\concat\sigma_p^*)$.
Conjunction forces both halves valid on the same $m^*$.
\emph{New message, both halves fresh}: feeding the output to $B_c$ (which runs
$A$ and simulates $\mathsf{Sign}_p$ honestly) converts $A$'s win into an ES256
forgery, up to a hash-collision term $q_h\cdot 2^{-\sigma/2}$.
\emph{Reused half}: accepting a reused classical half on a new message fails
$\mathsf{Ver}_p$ and vice versa. \emph{Malleability}: ECDSA admits
$(r,s)\mapsto(r,-s)$ and Falcon has encoding quirks, but modifying one block
leaves the other bound to $m$; the result is at most an existing transcript,
never a \emph{new-message} forgery. Hence the win event is contained in the
intersection of both component events and, for the conjunction rule,
$\varepsilon \le \min(\ldots) \le$ weaker component; substituting a full break
of one side leaves the bound at the other side's advantage. The composite
never advertises more security than ES256 (we claim EUF-CMA only, not SUF).
\end{proof}

\section{Channel binding: closing the bearer-replay window}
\begin{theorem}[OIDC over hybrid TLS inherits Theorem~\ref{thm:hybrid}]
\label{thm:bind}
If the token's signed input includes a TLS exporter-derived channel-binding
value $cb = \mathsf{KDF}(e, \text{label}, \text{ctx})$, then a token recorded
on one TLS connection fails verification on any other connection, and
transport secrecy of the flow reduces to Theorem~\ref{thm:hybrid} under the
TLS 1.3 PRF idealization with $K = \mathsf{HKDF}$-$\mathsf{Extract}(ss_h)$.
\end{theorem}

\begin{proof}[Proof sketch]
TLS's transcript hash binds the hybrid key share, so E3 cannot downgrade to a
classical-only share (distinct length, required by both peers). The JWT layer
without $cb$ has a compositional gap: a stolen $\mathit{id\_token}$ is
replayable to any RP that does not check $\mathit{c\_hash}/\mathit{at\_hash}$.
Including $cb$ in the signed input makes verification succeed only under the
exporter value captured on the signing connection; a replay carries a
different $cb$ and fails $\mathsf{Ver}_H$ by Theorem~\ref{thm:forge}.
\end{proof}

\section{HMK-1: the constant-time master-key wrap}
\label{sec:hmk}
HMK-1 seals a 32-byte master key under the hybrid secret
$ss = ss_c \concat ss_p \in \{0,1\}^{64}$ produced by Alg.~3. It is a
\emph{new} composition; no pre-existing block or stream cipher is used to
encrypt plaintext.

\subsection{Number theory: the $\F_p$ combiner}
Let $p = 2^{256} - 2^{32} + 977$ (NIST P-256 field prime) and let
$N: \{0,1\}^{256} \to \F_p^{\times}$ be reduction mod $p$ followed by the
branchless normalization $0 \mapsto 1$.

\begin{definition}[Field combiner]
$K_f = N(ss_c) \cdot N(ss_p) \bmod p$, computed by CIOS Montgomery
multiplication with runtime constants
$n_0^{-1} = 0\texttt{x}895950202ddacacf$ and $R^2 \bmod p$.
\end{definition}

\begin{proposition}[One-break property]
\label{prop:onebreak}
If exactly one input is known to the adversary and the other is uniform in
$\{0,1\}^{256}$, then $K_f$ is uniform over $\F_p^{\times}$ up to bias
$O(1/p)$. If both are unknown and independent, $K_f$ is uniform over
$\F_p^{\times}$ up to the same bias.
\end{proposition}

\begin{proof}
$N$ maps onto $\F_p$ with at most two preimages per output (bias $\le
2^{32}/2^{256}$) and $\{0,1\}^{256}\setminus N^{-1}(0)$ onto
$\F_p^{\times}$ after normalization. For fixed $a \ne 0$ the map
$b \mapsto ab \bmod p$ is a bijection of $\F_p$ (multiplicative inverse
exists), so a uniform $b$ maps to a uniform $K_f$ regardless of $a$ being
known. With both unknown, for every $y \ne 0$ there are exactly $p-1$ pairs
$(a,b)$ with $ab = y$ and $2p-1$ pairs with $ab = 0$; after normalization
this is uniform minus $O(1/p)$ mass on $\{1\}$.
\end{proof}

\begin{remark}
Unlike XOR-combining, the field product gives a one-break argument that is
one line of number theory --- and it is computable in constant time with a
fixed 4-limb loop and branchless reduction (Sec.~\ref{sec:sidechannel}).
\end{remark}

\subsection{Linear algebra: the $\F_2$ mixing layer}
\begin{definition}[Mixing matrix]
\label{def:matrix}
View the 64-byte secret as $x \in \F_2^{512}$ (eight 64-bit words, little
endian). Let
$T = I + S_1R^{13} + S_2R^{11} + S_4R^{31} + S_5R^{47}$, where $S_k$ cyclically
permutes words and $R^r$ rotates bits within a word; both are permutation
matrices over $\F_2$. The mixing layer is $M = T^3$ and $u = Mx$.
\end{definition}

\noindent\textbf{Verified properties} (exactly, by Gaussian elimination over
$\F_2$ and exhaustive basis-image enumeration --- \texttt{tools/linspread\_check.py},
re-checked in the C self-test):
\begin{itemize}
  \item $\mathrm{rank}(M) = 512$: $M$ is invertible (bijective diffusion).
  \item Every output word depends on all eight input words.
  \item A single input-bit flip changes at least 17 of the 64 output bytes
        (measured minimum over all 512 basis vectors).
  \item Cost: 3 rounds $\times$ 5 XOR-rotates per word $=$ 120 straight-line
        operations --- no branches, no tables.
\end{itemize}

\subsection{Construction}
\begin{algorithm}[t]
\caption{HMK-1: \texttt{wrap} / \texttt{unwrap} (fixed lengths:
$ss{=}64$, $mk{=}32$, $n{=}16$, $tag{=}32$ bytes)}
\label{alg:wrap}
\begin{algorithmic}[1]
\STATE $K_f \gets N(ss_c)\cdot N(ss_p) \bmod p$ \COMMENT{F\textsubscript{p} combiner}
\STATE $u \gets M\cdot ss$ \COMMENT{rank-512 F\textsubscript{2} mix}
\STATE $RK \gets \mathsf{HMAC}(K_f, u \| \texttt{"hmk/v1/rk"} \| i),\ i \in \{0,1\}$
\STATE $(kek, kmac) \gets RK_0 \| RK_1$
\STATE $KS \gets \mathsf{HMAC}(kek, \texttt{"hmk/v1/ks"} \| n)$ \COMMENT{1 block}
\STATE $ct \gets mk \oplus KS$
\STATE $tag \gets \mathsf{HMAC}(kmac, \texttt{"hmk/v1/tag"} \| n \| ct \| aad)$
\STATE \textbf{return} $(ct, tag)$
\STATESK $\hat{mk} \gets ct \oplus KS$;\quad
        $c \gets \mathsf{ct\_memcmp}(tag, \mathsf{HMAC}(kmac, \ldots))$
\STATESK $mk_{out} \gets 0^{32}$;\quad
        $\mathsf{cmov}(mk_{out}, \hat{mk}, [c = 0])$;\quad \textbf{return} $c$
\end{algorithmic}
\end{algorithm}

\noindent The AAD binds the record to a context (e.g.\ the issuer identity);
the nonce must be fresh per wrap under the same secret (standard one-time-pad
discipline), and both are covered by the tag.

\section{Side-channel analysis of the C implementation}
\label{sec:sidechannel}
The Cryptographer track asks for an algorithm that \emph{prevents} side-channel
attacks; we treat this as an engineering claim that must be measured.

\begin{table}[h]
\centering
\caption{Implementation-level vulnerability windows and their closures in
\texttt{csrc/}.}
\label{tab:sidechannel}
\begin{tabular}{@{}p{5.4cm}p{4.6cm}p{5.2cm}@{}}
\toprule
\textbf{Window} & \textbf{Naive behavior} & \textbf{HMK-1 closure} \\
\midrule
Short-circuit tag comparison &
  early exit on first mismatch; latency grows with matched prefix &
  \texttt{hmk\_ct\_memcmp}: single accumulator, full scan, zero exits \\
Secret-dependent branches &
  \texttt{if (x < p)} reductions, early returns &
  mask/cmov arithmetic; one return point per function \\
Secret-indexed tables (cache timing) &
  log/exp tables in field arithmetic &
  CIOS Montgomery: fixed 4-limb loops, no tables \\
Ordering oracle (verify-vs-derive) &
  reject before doing expensive work &
  identical statement path on every call; derive $\to$ keystream $\to$ compare \\
Partial key disclosure on failure &
  candidate left in memory &
  \texttt{cmov} from zeros: key published only on success, else zeroized \\
Variable-length encodings &
  parse/timing differentials &
  fixed-length asserts at every boundary (Go and C) \\
Compiler elision of wipes &
  \texttt{memset} removed as dead &
  \texttt{volatile} pointer wipe (\texttt{hmk\_secure\_zero}) \\
\bottomrule
\end{tabular}
\end{table}

\subsection{Measured results}
Reference machine: x86-64, GCC 13, \texttt{-O2}; medians of
$400 \times 512$ calls (\texttt{csrc/timing\_demo.c}).

\begin{table}[h]
\centering
\caption{Tag rejection latency vs.\ matching prefix length $k$: the naive
comparator is a byte-by-byte oracle; HMK-1 is flat.}
\label{tab:prefix}
\begin{tabular}{@{}lcccccccc@{}}
\toprule
$k$ & 0 & 1 & 2 & 4 & 8 & 16 & 31 & \\
\midrule
naive early-exit (ns) & 2 & 2 & 3 & 4 & 7 & 12 & 22 & $\nearrow$ linear \\
HMK-1 \texttt{ct\_memcmp} (ns) & 1 & 1 & 1 & 1 & 1 & 1 & 1 & flat \\
\bottomrule
\end{tabular}
\end{table}

\begin{table}[h]
\centering
\caption{Full \texttt{hmk\_unwrap} latency: success and every failure mode
execute the same path (400 interleaved rounds $\times$ 512 calls).}
\label{tab:unwrap}
\begin{tabular}{@{}lccccc@{}}
\toprule
case & correct tag & mismatch @0 & mismatch @31 & wrong secret & spread \\
\midrule
median latency & 8.57\,$\mu$s & 8.57\,$\mu$s & 8.58\,$\mu$s & 8.57\,$\mu$s &
  $0.2\%$ \\
\bottomrule
\end{tabular}
\end{table}

\begin{table}[h]
\centering
\caption{TVLA-style Welch $t$-test over the interleaved samples (threshold
$|t| < 4.5$, as in first-order side-channel certification).}
\label{tab:tvla}
\begin{tabular}{@{}lccc@{}}
\toprule
comparison & $|t|$ & threshold & verdict \\
\midrule
correct vs mismatch@0  & 1.25 & 4.5 & PASS \\
correct vs mismatch@31 & 0.53 & 4.5 & PASS \\
correct vs wrong secret & 1.06 & 4.5 & PASS \\
\bottomrule
\end{tabular}
\end{table}

Sampling uses the standard TVLA discipline: all cases are measured
\emph{interleaved in rotating order} so CPU-frequency drift biases every
distribution equally (a fixed measurement order falsely flagged
$|t| = 8.96$ during development — methodology, not code, was the cause).
The naive comparator in the same harness leaks a monotone slope that grows
$6$--$11\times$ from $k{=}0$ to $k{=}31$ ($2\to22$\,ns on the reference
run) — an oracle for byte-by-byte tag recovery — while HMK-1's rejection
latency is flat at 1\,ns for every prefix position.

The Go test \texttt{TestHMKUnwrapConstantTime} re-measures medians through
cgo on every \texttt{go test} run and fails if the spread exceeds $2\times$
(observed spread in CI-style runs: $0.1\%$). A deliberately vulnerable
comparator is shipped alongside (\texttt{naive\_verify}) so reviewers can
confirm the harness \emph{can} detect a leak.

\section{Evaluation: claim $\to$ proof $\to$ test}
\label{sec:trace}
\begin{table}[h]
\centering
\footnotesize
\setlength{\tabcolsep}{3pt}
\caption{Traceability of every claim to a proof and a named test.}
\label{tab:trace}
\begin{tabular}{@{}llp{7.4cm}@{}}
\toprule
\textbf{Claim} & \textbf{Proof} & \textbf{Test} \\
\midrule
E1: pq-break $\Rightarrow$ sessions safe & Thm~\ref{thm:hybrid}, Cor.~\ref{corr:e1} & \path{TestE1ClassicalHoldsUnderPQBreak} \\
E2: classical-break $\Rightarrow$ safe & Thm~\ref{thm:hybrid}, Cor.~\ref{corr:e2} & \path{TestE2HybridKeySurvivesClassicalBreak} \\
G2: tamper $\Rightarrow \bot$, fixed lengths & Thm~\ref{thm:hybrid} Games 0--3 & \path{TestCompositeKEMTamperIndist}, \path{TestFixedLengthEncodings} \\
G3: conjunction acceptance & Alg.~4 & \path{TestHybridVerifyRejectsEitherHalfMissing} \\
G4: weakest link & Thm~\ref{thm:forge} & \path{TestSignatureForgeryRequiresFreshECDSA} \\
Replay rejection & Thm~\ref{thm:bind} & \path{TestChannelBoundJWTRejectsReplay} \\
Wrap round-trip \& rejection surface & Alg.~\ref{alg:wrap} & \path{TestHMKWrapUnwrap}, C self-test (21 checks) \\
Alg.\ 3 end-to-end master-key recovery & Sec.~\ref{sec:hmk} & \path{TestHMKEndToEndWithHybridKEM} \\
Constant-time unwrap & Table~\ref{tab:sidechannel} & \path{TestHMKUnwrapConstantTime}, \texttt{hmk\_timing} \\
$M$ rank 512, full support & Def.~\ref{def:matrix} & \path{linspread\_check.py} + C self-test \\
F\textsubscript{p} arithmetic correctness & Prop.~\ref{prop:onebreak} & Python KAT + RFC 4231 vectors \\
\bottomrule
\end{tabular}
\end{table}

\section{What this work does not claim}
\begin{itemize}
  \item No security if \emph{both} KEM halves break; no SUF-CMA or one-more
        unforgeability (EUF-CMA at the weakest link only).
  \item The reference implementation realizes components as idealized RO
        stand-ins with break flags; concrete ML-KEM/ML-DSA integration
        changes no algorithm or proof but is not wired into a live OIDC
        provider in this submission.
  \item HMK-1's constant-time claim covers control flow and memory-access
        patterns in the C source as compiled by GCC/Clang at \texttt{-O2};
        it is not a formally verified assembly guarantee, and cross-border
        power/EM analysis is out of scope.
  \item The browser demo computes the same construction with BigInt/WebCrypto
        for visualization; it is not constant-time (only the C core carries
        that claim).
\end{itemize}

\section{Reproducibility}
\begin{lstlisting}[caption={Repository layout and commands},label={repo}]
csrc/            constant-time C core: make test && make timing
pkg/hybridoidc/  Alg 1-4 reference impl + cgo bindings: go test ./...
cmd/hmktui/      interactive TUI demo:        go run ./cmd/hmktui
site/            demo website (static, Vercel-ready)
docs/            full formal security specification (Markdown)
paper/           this paper:                  tectonic main.tex
tools/           independent vector/matrix checkers (Python, Node)
\end{lstlisting}

\begin{thebibliography}{9}
\bibitem{bindel2018}
S.~Bindel, J.~Brendel, M.~Fischlin, B.~González, U.~Maurer,
\emph{Hybrid Key Encapsulation Mechanisms and Authenticated Key Exchange},
PQCrypto 2019.

\bibitem{iethybrid}
IETF, \emph{Hybrid key exchange in TLS 1.3},
draft-ietf-tls-hybrid-design.

\bibitem{oidc}
OpenID Foundation, \emph{OpenID Connect Core 1.0}, 2014.

\bibitem{rfc8446}
T.~Rescorla, \emph{The Transport Layer Security (TLS) Protocol Version 1.3},
RFC 8446, 2018.

\bibitem{rfc9266}
N.~Sullivan, \emph{Channel Bindings for TLS 1.3}, RFC 9266, 2022.

\bibitem{fips180}
NIST, \emph{Secure Hash Standard (SHA-256)}, FIPS 180-4.

\bibitem{rfc2104}
H.~Krawczyk, M.~Bellare, R.~Canetti, \emph{HMAC: Keyed-Hashing for Message
Authentication}, RFC 2104, 1997.

\bibitem{chen2025}
Y.~Chen, S.~Thipireddy, A.~Ye, \emph{Hybrid Post-Quantum Variation of
OpenID}, MIT 6.5610, May 2025 (companion paper whose Alg.\ 3/4 we prove and
implement).

\bibitem{kos}
{\c{C}}.~Ko{\c{c}}, C.~Paar \emph{How to Do Montgomery Multiplication},
CHES 1998 (CIOS schedule).
\end{thebibliography}

\end{document}
